% ====================================================================
\section{Local modules: storage, transport, logic and reading}\label{sec:modules}
% ====================================================================

Every statement in this section concerns a forced evolution (\cref{def:controlled}). Every
contract holds for \emph{every} history of the ports that it declares arbitrary, so contracts can be
composed through \cref{lem:locality}.

\subsection{Cells, junctions and modes}

\begin{definition}[\textsf{EDGE} cell]\label{def:cell}
A \emph{cell} $X$ is a pair of adjacent vertices, its \emph{endpoints}. Each endpoint has exactly
four neighbours: its partner, one \emph{data port} outside the cell, and two private control slots.
If $Y$ is the cell, junction or external vertex attached to an endpoint of $X$, we write $\ep XY$
for that endpoint. Its data port is the vertex of $Y$ adjacent to it, written $\ep YX$ when $Y$ is a
cell. The cell is \emph{uniform with value $x$}, written $X=xx$, if both endpoints are in state $x$.
When the two endpoints differ we list them in the order stated in each proof.
\end{definition}

Both endpoints of a cell always receive the same pair of control bits. Three pairs are used; they
are called \emph{modes}:
\[
\HOLD=(0,1),\qquad \RESET=(0,0),\qquad \LOAD=(1,1),
\]
and we write $\FAV0=\RESET$ and $\FAV1=\LOAD$.
The pair is \emph{ordered}: in \HOLD\ the first slot of each endpoint carries $0$ and the second
carries $1$. The \emph{mode word} of $X$ is the map $\mu_X$ from control times to modes.

\begin{lemma}[Modes]\label{lem:modes}
Let $v$ be an endpoint of a cell, with partner $v'$ and data port $u$. At control time $t$,
\[
x_v(t+1)=\begin{cases}
\AND_3\bigl(x_v(t),x_{v'}(t),x_u(t)\bigr)&\text{in }\RESET,\\
\MAJ_3\bigl(x_v(t),x_{v'}(t),x_u(t)\bigr)&\text{in }\HOLD,\\
\OR_3\bigl(x_v(t),x_{v'}(t),x_u(t)\bigr)&\text{in }\LOAD.
\end{cases}
\]
\end{lemma}

\begin{proof}
By \eqref{eq:maj5}, $x_v(t+1)$ is the majority of five bits: three arguments and the two control
bits. With $k$ control bits equal to one, the result is one exactly when at least $3-k$ of the three
arguments are one.
\end{proof}

\begin{definition}[Junction]\label{def:junction}
A \emph{junction} $g$ is a single vertex adjacent to one endpoint of each of three cells, its
\emph{arms}, and to one control slot $h_g$. By \eqref{eq:maj5},
\[
x_g(t+1)=\MAJ_5\bigl(x_g(t),a_1(t),a_2(t),a_3(t),h_g(t)\bigr),
\]
where $a_1,a_2,a_3$ are the arm
endpoints adjacent to $g$.
\end{definition}

A \emph{cell network} is a controlled tree built from cells and junctions. Its control slots are
the two slots of each cell endpoint and the slot $h_g$ of each junction. Some cell endpoints have,
as data port, an \emph{external port}: a vertex whose history is supplied from outside. In the
processor these are the far clocks of \cref{sec:processor}. A \emph{schedule} of a cell network
consists of a mode word for each cell, a word for each junction control, and a history for each
external port.

\begin{lemma}[Mode-level duality]\label{lem:duality}
Let $x(\cdot)$ be the forced evolution of a cell network under a schedule with mode words $\mu_X$,
junction words $h_g$ and external port histories $y_u$. Define $\mu^\ast_X$ by exchanging $\RESET$
and $\LOAD$ in $\mu_X$ (\HOLD\ is unchanged), $h^\ast_g=1-h_g$, and $y^\ast_u=1-y_u$. Then
$\mathbf 1-x(\cdot)$ is the forced evolution under $(\mu^\ast,h^\ast,y^\ast)$ from the initial
configuration $\mathbf 1-x(0)$.
\end{lemma}

\begin{proof}
By \cref{lem:modes} an endpoint applies $f_\mu\in\{\AND_3,\MAJ_3,\OR_3\}$. De Morgan's laws give
$1-\AND_3(\bar a)=\OR_3(a)$ and $1-\OR_3(\bar a)=\AND_3(a)$, and $\MAJ_3$ is self-dual. Hence
$1-f_\mu(\mathbf 1-a)=f_{\mu^\ast}(a)$. At junctions, $\MAJ_5$ is self-dual. Induction on time
concludes.
\end{proof}

\begin{remark}\label{rem:dualcontrols}
\Cref{lem:duality} is a statement about \emph{modes}, not about the control bits themselves.
Literally complementing the control bits of a cell would turn $\HOLD=(0,1)$ into $(1,0)$. That pair
is also a hold mode, but it is a different pair of slot words. We never do this: in every schedule,
\HOLD\ is realised by $(0,1)$. The effect of duality on admissible timings is computed in
\cref{sec:phasecalc}.
\end{remark}

For a \emph{background} $c\in\{0,1\}$ we write $b=1-c$ throughout. All primitives below are stated for
an arbitrary background $c$, using $\FAV c$, $\FAV b$ and junction controls equal to $b$ or $c$. By
\cref{lem:duality}, the instance with $c=1$ is exactly the dual of the instance with $c=0$, so every
proof below is written for $c=0$ only.

\begin{lemma}[Frame]\label{lem:frame}
Let $t$ be a control time.
\begin{enumerate}[label=(\alph*)]
\item\label{frame:hold} A uniform cell $X=xx$ in \HOLD\ at control time $t$ satisfies $X=xx$ at
  time $t+1$, whatever its two data ports are.
\item\label{frame:junction} If a junction satisfies $x_g(t)=h_g(t)=c$, and at most two of its three
  arm endpoints are in state $b$ at time $t$, then $x_g(t+1)=c$.
\end{enumerate}
\end{lemma}

\begin{proof}
\ref{frame:hold} By \cref{lem:modes}, each endpoint computes $\MAJ_3(x,x,\cdot)=x$.
\ref{frame:junction} At least three of the five arguments of $\MAJ_5$ equal $c$.
\end{proof}

In every later proof, a cell that a schedule does not mention is in \HOLD. ``By the frame lemma''
means an application of \cref{lem:frame}. Only the arms of a junction and its control can change
its state, so to keep an untouched junction at $c$ it suffices to check, at every control time, that
at most two of its arm endpoints are $b$ and that its control is at rest at $c$.

\subsection{Storage}

\begin{lemma}[Forcing]\label{lem:forcing}
Let $X$ be in $\FAV z$ at the control times $[t,t+n)$ with $n\ge 2$, and suppose that the data port
of some endpoint $v$ of $X$ equals $z$ at control time $t$. Then $X=zz$ at every time in
$[t+2,t+n]$, whatever the initial state of $X$ and whatever its other data port does. If moreover
$X=zz$ at time $t$, then $X=zz$ at every time in $[t,t+n]$.
\end{lemma}

\begin{proof}
For $z=0$: $x_v(t+1)=\AND_3(\cdot,\cdot,0)=0$. Then $x_{v'}(t+2)=\AND_3(\cdot,x_v(t+1),\cdot)=0$
and $x_v(t+2)=\AND_3(x_v(t+1),\cdot,\cdot)=0$. Once both endpoints are $0$, $\AND_3$ with self
argument $0$ keeps them at $0$. The case $z=1$ is dual.
\end{proof}

\begin{definition}[\prim{SET}]\label{def:set}
$\prim{SET}(X\mathbin{:=}z)$ puts $X$ in $\FAV z$ for three consecutive control times, while some data
port of $X$ is known to equal $z$ at the first of them. By \cref{lem:forcing}, two transitions
suffice and the third preserves $X=zz$. Three transitions are used so that the timing matches the
phase calculus of \cref{sec:phasecalc}.
\end{definition}

\begin{remark}\label{rem:forcing-hyp}
The original manuscript states the forcing hypothesis as ``the port stays $z$ for two transitions''.
\Cref{lem:forcing} needs the port value at the first control time only. This is a weaker hypothesis,
so every application in the original remains valid.
\end{remark}

\begin{lemma}[Copy]\label{lem:copy}
Let $X=cc$ be in $\FAV b$ during $[t,t+n)$ with $n\ge 2$. Let the data port of one endpoint
$v$ be constantly $x$ during $[t,t+n)$, and the data port of the other endpoint $v'$ be constantly
$c$. Then $X=xx$ at every time in $[t+2,t+n]$.
\end{lemma}

\begin{proof}
For $c=0$ the mode is $\OR_3$. Then $x_v(t+1)=\OR(0,0,x)=x$ and $x_{v'}(t+1)=\OR(0,0,0)=0$.
Next, $x_v(t+2)=\OR(x,0,x)=x$ and $x_{v'}(t+2)=\OR(0,x,0)=x$. Afterwards
$x_v=\OR(x,x,x)=x$ and $x_{v'}=\OR(x,x,0)=x$.
\end{proof}

$\prim{COPY}$ into $X$ is \cref{lem:copy} with $n=3$. Its dual form starts from $X=bb$, uses
$\FAV c$, has its second port constantly $b$, and again yields $X=xx$. The source of a copy is never
modified by it: it is either an external port or a uniform cell in \HOLD.

\subsection{Transport}

\begin{lemma}[\prim{MOVE}]\label{lem:move}
Let $X,Y$ be cells joined by the edge $\ep XY\ep YX$. Let $\ep X{\mathrm f}$ and $\ep Y{\mathrm f}$ be
their other endpoints, whose data ports are $u_X$ and $u_Y$. Assume that at time $t$ we have
$X=xx$ and $Y=zz$, and that $u_Y=z$ at control times $t,t+1,t+2$. The history of $u_X$ is arbitrary.
Put $X$ in $\FAV z$ and $Y$ in $\FAV{1-z}$ during $[t,t+3)$. Then at time $t+3$ we have $X=zz$ and
$Y=xx$. This is $\prim{MOVE}_z(X\to Y)$.
\end{lemma}

\begin{proof}
For $z=0$, writing $P(t)$ for the history of $u_X$, the states are:
\begin{center}\small
\begin{tabular}{ccccc}\toprule
state time & $\ep X{\mathrm f}$ & $\ep XY$ & $\ep YX$ & $\ep Y{\mathrm f}$\\\midrule
$t$   & $x$ & $x$ & $0$ & $0$\\
$t+1$ & $x\wedge P(t)$ & $0$ & $x$ & $0$\\
$t+2$ & $0$ & $0$ & $x$ & $x$\\
$t+3$ & $0$ & $0$ & $x$ & $x$\\\bottomrule
\end{tabular}
\end{center}
Transition $t$: $\ep XY=\AND(x,x,0)=0$ reads the old $\ep YX=0$, while $\ep YX=\OR(0,0,x)=x$ reads the
old $\ep XY=x$; and $\ep Y{\mathrm f}=\OR(0,0,u_Y)=0$.
Transition $t+1$: $\ep X{\mathrm f}=\AND(\cdot,0,\cdot)=0$; $\ep YX=\OR(x,0,0)=x$;
$\ep Y{\mathrm f}=\OR(0,x,u_Y)=x$ because $u_Y=0$.
Transition $t+2$: $X$ stays $00$ under $\AND$. On $Y$, $\OR(x,x,\cdot)=x$ when $x=1$, and every
argument is $0$ when $x=0$.
\end{proof}

\begin{remark}[No pause inside a move]\label{rem:nopause}
The three transitions of a move must be consecutive. In the case $P(t)=x$, the state after the
first transition is $(\ep X{\mathrm f},\ep XY,\ep YX,\ep Y{\mathrm f})=(x,0,x,0)$. If both cells were then
put in \HOLD\ for one step, \cref{lem:modes} would give
\[
\bigl(\MAJ_3(x,0,x),\MAJ_3(0,x,x),\MAJ_3(x,0,0),\MAJ_3(0,x,0)\bigr)=(x,x,0,0),
\]
which undoes the transport. Alignment waits (\cref{lem:waits}) are therefore inserted only between
primitives, at configurations proved to be stable.
\end{remark}

\begin{lemma}[\prim{ROUTE}]\label{lem:route}
Let $g$ be a junction with arms $A$ (source), $B$ (destination) and $C$ (bias). Assume at time $s$:
\[
x_g=c,\qquad A=xx,\qquad B=cc,\qquad C=bb.
\]
Assume also that the data port of $\ep B{\mathrm f}$, the endpoint of $B$ away from $g$, equals $c$ at
control times $s+3,s+4,s+5$. The port of $\ep A{\mathrm f}$ is arbitrary. Apply:
\begin{center}\small
\begin{tabular}{lll}\toprule
control times & $h_g$ & modes\\\midrule
$[s,s+3)$ & $b$ & $A$ in $\FAV c$ (receive and erase)\\
$[s+3,s+6)$ & $c$ & $B$ in $\FAV b$ (send)\\\bottomrule
\end{tabular}
\end{center}
with all other cells in \HOLD. Then at time $s+6$ we have $x_g=c$, $A=cc$, $B=xx$ and $C=bb$. This is
$\prim{ROUTE}(g;A\to B;C)$, of duration $6$.
\end{lemma}

\begin{proof}
For $c=0$, with $P$ the arbitrary port of $\ep A{\mathrm f}$:
\begin{center}\small
\begin{tabular}{cccccc}\toprule
state time & $x_g$ & $\ep A{\mathrm f}$ & $\ep Ag$ & $\ep Bg$ & $\ep B{\mathrm f}$\\\midrule
$s$   & $0$ & $x$ & $x$ & $0$ & $0$\\
$s+1$ & $x$ & $x\wedge P$ & $0$ & $0$ & $0$\\
$s+2$ & $x$ & $0$ & $0$ & $0$ & $0$\\
$s+3$ & $x$ & $0$ & $0$ & $0$ & $0$\\
$s+4$ & $0$ & $0$ & $0$ & $x$ & $0$\\
$s+5$ & $0$ & $0$ & $0$ & $x$ & $x$\\
$s+6$ & $0$ & $0$ & $0$ & $x$ & $x$\\\bottomrule
\end{tabular}
\end{center}
$C=11$ throughout, by \cref{lem:frame}\ref{frame:hold}.
Receive, with $h_g=1$ and $C=1$: then $x_g(t+1)=\MAJ_5(x_g,\ep Ag,\ep Bg,1,1)$, which is $\OR_3$ of
$x_g$, $\ep Ag$ and $\ep Bg$. At control time $s$ this gives $x_g=\OR(0,x,0)=x$, while
$\ep Ag=\AND(x,x,x_g=0)=0$. At $s+1$, $x_g=\OR(x,0,0)=x$ and $\ep A{\mathrm f}=\AND(\cdot,0,\cdot)=0$.
At $s+2$ nothing changes. $B=00$ stays in \HOLD\ by the frame lemma, although its port $x_g$ may be
$1$.
Send, with $h_g=0$: the only possible one-votes at $g$ are $x_g$, $\ep Bg$ and $C$. At $s+3$ there are
at most two of them ($x_g=x$ and $C$), so $x_g$ becomes $0$, while $\ep Bg=\OR(0,0,x_g=x)=x$ and
$\ep B{\mathrm f}=\OR(0,0,0)=0$. At $s+4$ and $s+5$ the one-votes at $g$ are at most $\ep Bg$ and $C$,
since $x_g=0$, so $g$ stays $0$; and $B$ becomes and stays $xx$ as in the move.
\end{proof}

Note that during a route the junction itself carries the bit $x$ at times $s+1,s+2,s+3$. Its arms
are the only cells that read it.

The following lemma composes routes along a path of junctions. Its only external hypothesis concerns
the final destination; the intermediate far ports are handled inside the proof by the frame lemma.

\begin{lemma}[Transport along a chain]\label{lem:chain}
Let $g_1,\dots,g_n$ be junctions and $X_0,\dots,X_n$ cells such that, for $1\le\ell\le n$, the cells
$X_{\ell-1}$ and $X_\ell$ are arms of $g_\ell$, with third arm $Y_\ell$. For $\ell<n$, the endpoint
of $X_\ell$ away from $g_\ell$ is adjacent to $g_{\ell+1}$. Assume at time $s$:
\[
X_0=xx,\quad X_1=\dots=X_n=cc,\quad Y_1=\dots=Y_n=bb,\quad x_{g_1}=\dots=x_{g_n}=c,
\]
and that the data port of the endpoint of $X_n$ away from $g_n$ is $c$ during
$[s+6n-3,s+6n)$. Apply $\prim{ROUTE}(g_\ell;X_{\ell-1}\to X_\ell;Y_\ell)$ during
$[s+6(\ell-1),s+6\ell)$ for $\ell=1,\dots,n$, with $h_{g_\ell}=c$ outside its own route. Then at
time $s+6n$ we have $X_n=xx$, $X_0=\dots=X_{n-1}=cc$, all $g_\ell=c$, and the $Y_\ell$ unchanged.
\end{lemma}

\begin{proof}
By induction on $\ell$ we prove the invariant $J_\ell$: at time $s+6\ell$, $X_\ell=xx$, all other
$X$'s are $cc$, all junctions are $c$, and the $Y$'s are unchanged. $J_0$ is the hypothesis. Assume
$J_{\ell-1}$, and consider control times $[s+6(\ell-1),s+6\ell)$. We verify the hypotheses of
\cref{lem:route} at $g_\ell$; its source $X_{\ell-1}$ may have any far port. Its destination
$X_\ell$ has far port $g_{\ell+1}$ when $\ell<n$. The arms of $g_{\ell+1}$ are $X_\ell$,
$X_{\ell+1}=cc$ and $Y_{\ell+1}=bb$, and its control is $c$. So at most two of its arm endpoints
are $b$, and the frame lemma keeps $g_{\ell+1}=c$ throughout. When $\ell=n$ the far port is $c$ by
hypothesis. Every other junction $g_{\ell'}$ keeps state $c$ by the frame lemma: for
$\ell'\ge\ell+2$ its arms are unchanged, $c$ and $b$; for $\ell'=\ell-1$ its arms are $X_{\ell-2}=cc$,
$Y_{\ell-1}=bb$ and $X_{\ell-1}$, which is $x$ or $c$. Every other cell is a uniform cell in \HOLD.
\Cref{lem:route} gives $J_\ell$.
\end{proof}

The lemma does not control the neighbour of $X_0$ away from $g_1$. That vertex sees $X_0$ change from
$x$ to $c$, and every application checks it separately.

\subsection{Logic}

\begin{lemma}[\prim{GATE}]\label{lem:gate}
Let $g$ be a junction with arms $A$, $B$ (operands) and $C$ (output). Assume at time $s$:
\[
x_g=c,\qquad A=xx,\qquad B=yy,\qquad C=cc,
\]
and that the data port of $\ep C{\mathrm f}$ equals $c$ during $[s,s+5)$. The operands' far ports are
arbitrary. Apply:
\begin{center}\small
\begin{tabular}{llll}\toprule
control times & $C$ & $A,B$ & $h_g$\\\midrule
$s$ & $\FAV b$ (pre-enable) & \HOLD & $c$\\
$[s+1,s+4)$ & $\FAV b$ & $\FAV c$ (core) & $b$\\
$s+4$ & $\FAV b$ & \HOLD & $c$\\\bottomrule
\end{tabular}
\end{center}
Then at time $s+5$ we have $x_g=c$, $A=B=cc$ and $C=dd$, where $d=x\wedge y$ if $c=0$ and
$d=x\vee y$ if $c=1$. This is $\prim{GATE}_c(g;A,B\to C)$, of duration $5$.
\end{lemma}

\begin{proof}
For $c=0$, with $d=x\wedge y$ (the operand far endpoints are omitted; they reach $0$ by $s+3$):
\begin{center}\small
\begin{tabular}{cccccc}\toprule
state time & $x_g$ & $\ep Ag$ & $\ep Bg$ & $\ep Cg$ & $\ep C{\mathrm f}$\\\midrule
$s$   & $0$ & $x$ & $y$ & $0$ & $0$\\
$s+1$ & $0$ & $x$ & $y$ & $0$ & $0$\\
$s+2$ & $d$ & $0$ & $0$ & $0$ & $0$\\
$s+3$ & $0$ & $0$ & $0$ & $d$ & $0$\\
$s+4$ & $0$ & $0$ & $0$ & $d$ & $d$\\
$s+5$ & $0$ & $0$ & $0$ & $d$ & $d$\\\bottomrule
\end{tabular}
\end{center}
At $s$, $\MAJ_5(0,x,y,0,0)=0$, and $C$ in $\OR_3$ sees only zeros. At $s+1$, with $h_g=1$,
$\MAJ_5(0,x,y,0,1)=x\wedge y$, while each operand's near endpoint computes $\AND(\cdot,\cdot,x_g=0)=0$
and its far endpoint $\AND(\cdot,\cdot,P)$. At $s+2$, $\MAJ_5(d,0,0,0,1)=0$ because $d+1\le 2$; the
operands are now entirely $0$; and $\ep Cg=\OR(0,0,x_g=d)=d$. At $s+3$, $\MAJ_5(0,0,0,d,1)=0$ and
$\ep C{\mathrm f}=\OR(0,d,0)=d$. At $s+4$, $h_g=0$ and $\MAJ_5(0,0,0,d,0)=0$, while $C$ keeps $dd$
under $\OR$.
\end{proof}

For $c=1$ the dual gate computes $x\vee y$ on background one. Duality proves a dual
\emph{operation}. It adds no negation of an arbitrary input and does not change the majority rule.

\subsection{Reading with restoration of the isolating cells}

Reading a register must preserve its content and restore the cells that isolate it from the bus. A
naive read that leaves a zero-valued source connected to a zero-valued copy along an entirely zero
path cannot be undone later. The following operation restores the buffer before sending, and
restores the front cell during sending.

\begin{lemma}[\prim{READ}]\label{lem:read}
Let $g$ be a junction with arms $A$, $B$, $C$, and let the side branch continue $C\,$--$\,K\,$--$\,R$,
that is, $\ep CK$ is adjacent to $\ep KC$ and $\ep KR$ to $\ep RK$:
\[
A\;-\;g\;-\;B,\qquad g\;-\;C\;-\;K\;-\;R .
\]
Assume at time $s$:
\[
A=C=K=bb\ (\text{the pilot and the isolators}),\qquad B=cc,\qquad R=xx,
\]
with $x_g$ arbitrary. Assume the data port of $\ep A{\mathrm f}$ equals $c$ during $[s,s+14)$. The
data ports of $\ep B{\mathrm f}$ and $\ep R{\mathrm f}$ are arbitrary. $R$ is in \HOLD\ throughout.
Apply, with unmentioned cells in \HOLD:
\begin{center}\small
\begin{tabular}{lll}\toprule
control times & $h_g$ & modes\\\midrule
$s$ & $b$ & ---\\
$[s+1,s+4)$ & $b$ & $K$ in $\FAV c$\\
$[s+4,s+7)$ & $b$ & $K$ in $\FAV b$, $C$ in $\FAV c$\\
$s+7$ & $b$ & ---\\
$[s+8,s+11)$ & $b$ & $A$ in $\FAV c$\\
$[s+11,s+14)$ & $c$ & $A$ and $C$ in $\FAV b$\\\bottomrule
\end{tabular}
\end{center}
Then at time $s+14$ we have
\[
x_g=c,\qquad A=xx,\qquad B=cc,\qquad C=K=bb,\qquad R=xx.
\]
The bit $x$ is delivered to $A$, the source $R$ is preserved, and both isolators $C$ and $K$ are
restored; the pilot $A=bb$ is consumed. This is $\prim{READ}(g;A;C,K,R)$, of duration $14$.
\end{lemma}

\begin{proof}
For $c=0$ (so the pilot and isolators are $11$). Columns list $x_g$ and the endpoints; $B=00$ is
constant by the frame lemma, and so is $R=xx$.
\begin{center}\small
\begin{tabular}{cccccccc}\toprule
state time & $x_g$ & $\ep A{\mathrm f}$ & $\ep Ag$ & $\ep Cg$ & $\ep CK$ & $\ep KC$ & $\ep KR$\\\midrule
$s$    & any & $1$ & $1$ & $1$ & $1$ & $1$ & $1$\\
$s+1$  & $1$ & $1$ & $1$ & $1$ & $1$ & $1$ & $1$\\
$s+2$  & $1$ & $1$ & $1$ & $1$ & $1$ & $1$ & $x$\\
$s+3$  & $1$ & $1$ & $1$ & $1$ & $1$ & $x$ & $x$\\
$s+4$  & $1$ & $1$ & $1$ & $1$ & $1$ & $x$ & $x$\\
$s+5$  & $1$ & $1$ & $1$ & $1$ & $x$ & $1$ & $x$\\
$s+6$  & $1$ & $1$ & $1$ & $x$ & $x$ & $1$ & $1$\\
$s+7$  & $1$ & $1$ & $1$ & $x$ & $x$ & $1$ & $1$\\
$s+8$  & $1$ & $1$ & $1$ & $x$ & $x$ & $1$ & $1$\\
$s+9$  & $1$ & $0$ & $1$ & $x$ & $x$ & $1$ & $1$\\
$s+10$ & $1$ & $0$ & $0$ & $x$ & $x$ & $1$ & $1$\\
$s+11$ & $x$ & $0$ & $0$ & $x$ & $x$ & $1$ & $1$\\
$s+12$ & $0$ & $0$ & $x$ & $x$ & $1$ & $1$ & $1$\\
$s+13$ & $0$ & $x$ & $x$ & $1$ & $1$ & $1$ & $1$\\
$s+14$ & $0$ & $x$ & $x$ & $1$ & $1$ & $1$ & $1$\\\bottomrule
\end{tabular}
\end{center}
We justify each transition.
\begin{itemize}
\item \emph{Control time $s$.} With $\ep Ag=\ep Cg=h_g=1$, the junction becomes $1$ whatever its
  previous state.
\item \emph{Extraction, $[s+1,s+4)$, $K$ in $\RESET$.} The ports of $K$ are $\ep CK=1$ and
  $\ep RK=x$. At $s+1$, $\ep KR=\AND(1,1,x)=x$ and $\ep KC=\AND(1,1,1)=1$. At $s+2$,
  $\ep KC=\AND(1,x,1)=x$. At $s+3$, $K=xx$ is preserved. The junction stays $1$: its self, $\ep Ag$ and
  $h_g$ are three one-votes.
\item \emph{Transfer, $[s+4,s+7)$.} This is $\prim{MOVE}_1(K\to C)$ of \cref{lem:move}: source $K=xx$
  in \LOAD, destination $C=11$ in \RESET, and far destination port $x_g$, which stays $1$ because
  its self, $\ep Ag$ and $h_g$ are $1$. The source's far port is $\ep RK=x$. The result is $K=11$,
  $C=xx$, as listed.
\item \emph{Control time $s+7$.} All cells are in \HOLD\ and uniform; $x_g$ stays $1$.
\item \emph{Receive, $[s+8,s+11)$, $A$ in \RESET\ with far port $0$.} At $s+8$,
  $\ep A{\mathrm f}=\AND(1,1,0)=0$ and $\ep Ag=\AND(1,1,x_g=1)=1$. At $s+9$, $\ep Ag=\AND(1,0,1)=0$.
  Meanwhile $x_g$ stays $1$, because self and $h_g$ are $1$ and $\ep Ag(s+9)=1$. At $s+10$,
  $x_g(s+11)=\MAJ_5(1,0,0,x,1)=x$, since $\ep Ag=\ep Bg=0$ and $\ep Cg=x$.
\item \emph{Send and restore, $[s+11,s+14)$, $h_g=0$, $A$ and $C$ in \LOAD.} At $s+11$,
  $x_g=\MAJ_5(x,0,0,x,0)=0$; $\ep Ag=\OR(0,0,x_g=x)=x$; $\ep A{\mathrm f}=\OR(0,0,0)=0$;
  $\ep Cg=\OR(x,x,x)=x$; and $\ep CK=\OR(x,x,\ep KC=1)=1$. At $s+12$, the one-votes at $g$ are at
  most $\ep Ag$ and $\ep Cg$, so $x_g=0$; $\ep A{\mathrm f}=\OR(0,x,0)=x$; $\ep Ag=\OR(x,0,0)=x$; and
  $\ep Cg=\OR(x,1,0)=1$. At $s+13$, $\MAJ_5(0,x,0,1,0)=0$, and $A=xx$ and $C=11$ are preserved under
  $\OR$ with port values $0$.
\end{itemize}
$B=00$ is in \HOLD\ and uniform, so the bits $1$ or $x$ at $g$ do not affect it. $R$ is in \HOLD\ and
uniform, so $K$ does not affect it.
\end{proof}

\begin{remark}\label{rem:readwait}
The pilot, the isolators and the source are uniform cells in \HOLD, and the junction may start in
either state. So the read can be preceded by any number of all-\HOLD\ steps with $h_g=c$ and no
transient value of $g$ is needed. In the processor the read always starts with $x_g=c$.
\end{remark}
